\section{结果统计框架}

\subsection{评估原则}

本文在主结果中采用 5 折重复 10 次的 repeated stratified cross-validation 作为统一评估口径，共 50 个测试折，并报告均值、标准差与 95\% 置信区间。这样做的目标是降低单次切分带来的偶然性，让不同方法的比较建立在更稳定的统计基础上。

\subsection{统计量定义}

对于每个方法和每个指标，设 fold-level 结果为 $x_1, x_2, \dots, x_n$，则：

\begin{itemize}
    \item 均值：$\mu = \frac{1}{n} \sum_{i=1}^{n} x_i$
    \item 标准差：$\sigma = \sqrt{\frac{1}{n} \sum_{i=1}^{n} (x_i - \mu)^2}$
    \item 95\% CI：$\mu \pm 1.96 \cdot \frac{\sigma}{\sqrt{n}}$
\end{itemize}

其中 $n$ 为全部测试折数量（本文 $n=50$）。

\subsection{比较原则}

不同方法的结论采用以下顺序解释：

\begin{enumerate}
    \item 先比较 Accuracy、Balanced Accuracy 与 Macro-F1 的均值；
    \item 再看 95\% CI 是否明显重叠；
    \item 最后结合 fold-level 分布判断稳定性。
\end{enumerate}

若某方法仅在单次切分上表现更优，但在 repeated CV 下波动较大，则不将其作为主结论依据。

\subsection{显著性检验说明}

当需要比较相邻方法是否存在稳定差异时，我们采用基于 fold-level 指标的双侧置换检验（20,000 次置换）作为补充分析工具。该检验不依赖正态分布假设，适合小样本、多折结果场景。

\subsection{多指标统计结果}

\begin{table}[H]
\centering
\caption{四方法性能对比（5折×10次重复）}
\label{tab:results-main}
\begin{tabular}{lcccccc}
\toprule
方法 & Accuracy & 95\% CI & Balanced Accuracy & 95\% CI & Macro F1 & 95\% CI \\
\midrule
RNA-only & 0.8655 & [0.8530, 0.8780] & 0.8704 & [0.8570, 0.8838] & 0.8626 & [0.8487, 0.8765] \\
Concat & 0.9012 & [0.8906, 0.9118] & 0.9074 & [0.8962, 0.9186] & 0.8999 & [0.8882, 0.9116] \\
MOFA & 0.7857 & [0.7640, 0.8074] & 0.7889 & [0.7660, 0.8118] & 0.7659 & [0.7423, 0.7895] \\
Stacking & 0.8930 & [0.8814, 0.9046] & 0.8955 & [0.8832, 0.9078] & 0.8878 & [0.8749, 0.9007] \\
\bottomrule
\end{tabular}
\end{table}

从均值排序为：Concat > Stacking > RNA-only > MOFA。Concat 在三项指标均表现最佳。

\subsection{多指标置换检验}

\begin{table}[H]
\centering
\caption{四方法两两置换检验（多指标均值差）}
\label{tab:permtest-all}
\begin{tabular}{lccc}
\toprule
方法对比 & |ΔAccuracy| & |ΔBalancedAcc| & |ΔMacroF1| \\
\midrule
RNA vs Concat & 0.0357 & 0.0370 & 0.0373 \\
RNA vs MOFA & 0.0798 & 0.0815 & 0.0967 \\
RNA vs Stacking & 0.0275 & 0.0251 & 0.0252 \\
Concat vs MOFA & 0.1155 & 0.1185 & 0.1340 \\
Concat vs Stacking & 0.0082 & 0.0119 & 0.0121 \\
MOFA vs Stacking & 0.1073 & 0.1066 & 0.1219 \\
\bottomrule
\end{tabular}
\end{table}

从当前样本规模看，虽然 Concat 在均值上领先，但统计显著性尚不充分。我们据此把“稳定趋势”和“区间重叠关系”作为主结论依据。